\clearpage\phantomsection\label{couverture-programme}\pdfbookmark[0]{Repérage du programme}{bm-couverture-programme}\gdef\otChapNum{0}\pagestyle{otFront}
\otPageTitre{Repérage du programme}{Cinquième · Nouveau programme 2026}
Cette matrice relie les notions et les automatismes du programme aux
exercices du livre. Les cours complètent ces repères, notamment pour les
définitions, démonstrations et constructions.\par\bigskip
Les objectifs nouveaux sont intégrés : fractions de dénominateurs
quelconques, carrés et cubes, distributivité littérale, équations simples,
parallélisme et angles, hauteurs et médianes, cercle circonscrit, y compris
celui du triangle rectangle, caractérisations des parallélogrammes,
perspectives et patrons de prismes et de cylindres, moyenne simple, échelle
de probabilité, lecture de dépendances et répétitions inconditionnelles.\par\bigskip
\begin{attention}Les niveaux de difficulté et les défis sont des choix
d’entraînement. Ils ne constituent pas une liste d’obligations
supplémentaires au programme officiel.\end{attention}\clearpage
\otPageTitre{Repères et exercices}{Chapitres 1 à 3}\fontsize{10}{13}\selectfont
\otIntertitre{1. Opérations et priorités}\begin{tabularx}{\linewidth}{@{}Yl@{}}\toprule
Opérations et ordres de grandeur & \hyperref[exo-1-1]{1.1}, \hyperref[exo-1-2]{1.2}, \hyperref[exo-1-3]{1.3}, \hyperref[exo-1-11]{1.11}\\[2mm]
Division par un décimal & \hyperref[exo-1-4]{1.4}, \hyperref[exo-1-9]{1.9}, \hyperref[exo-1-12]{1.12}\\[2mm]
Priorités et vocabulaire & \hyperref[exo-1-5]{1.5}, \hyperref[exo-1-6]{1.6}, \hyperref[exo-1-7]{1.7}, \hyperref[exo-1-10]{1.10}\\[2mm]
Distributivité numérique et problèmes & \hyperref[exo-1-8]{1.8}, \hyperref[exo-1-13]{1.13}, \hyperref[exo-1-14]{1.14}, \hyperref[exo-1-15]{1.15}\\[2mm]
\bottomrule\end{tabularx}\par\vspace{4mm}
\otIntertitre{2. Multiples, diviseurs et partages}\begin{tabularx}{\linewidth}{@{}Yl@{}}\toprule
Division euclidienne et reste & \hyperref[exo-2-1]{2.1}, \hyperref[exo-2-5]{2.5}, \hyperref[exo-2-11]{2.11}\\[2mm]
Multiples, diviseurs et critères & \hyperref[exo-2-2]{2.2}, \hyperref[exo-2-3]{2.3}, \hyperref[exo-2-4]{2.4}, \hyperref[exo-2-6]{2.6}, \hyperref[exo-2-7]{2.7}\\[2mm]
Recherche organisée de solutions & \hyperref[exo-2-8]{2.8}, \hyperref[exo-2-9]{2.9}, \hyperref[exo-2-14]{2.14}\\[2mm]
Partages et phénomènes périodiques & \hyperref[exo-2-10]{2.10}, \hyperref[exo-2-12]{2.12}, \hyperref[exo-2-13]{2.13}, \hyperref[exo-2-15]{2.15}\\[2mm]
\bottomrule\end{tabularx}\par\vspace{4mm}
\otIntertitre{3. Découvrir les nombres relatifs}\begin{tabularx}{\linewidth}{@{}Yl@{}}\toprule
Droite graduée et références & \hyperref[exo-3-1]{3.1}, \hyperref[exo-3-2]{3.2}, \hyperref[exo-3-8]{3.8}\\[2mm]
Opposé et valeur absolue & \hyperref[exo-3-3]{3.3}, \hyperref[exo-3-4]{3.4}, \hyperref[exo-3-10]{3.10}\\[2mm]
Comparaison et encadrements & \hyperref[exo-3-5]{3.5}, \hyperref[exo-3-6]{3.6}, \hyperref[exo-3-7]{3.7}, \hyperref[exo-3-12]{3.12}, \hyperref[exo-3-13]{3.13}\\[2mm]
Signe et interprétation & \hyperref[exo-3-9]{3.9}, \hyperref[exo-3-11]{3.11}, \hyperref[exo-3-14]{3.14}, \hyperref[exo-3-15]{3.15}\\[2mm]
\bottomrule\end{tabularx}\par\vspace{4mm}
\clearpage
\otPageTitre{Repères et exercices}{Chapitres 4 à 6}\fontsize{10}{13}\selectfont
\otIntertitre{4. Additionner et soustraire des relatifs}\begin{tabularx}{\linewidth}{@{}Yl@{}}\toprule
Addition de relatifs & \hyperref[exo-4-1]{4.1}, \hyperref[exo-4-2]{4.2}, \hyperref[exo-4-3]{4.3}, \hyperref[exo-4-5]{4.5}, \hyperref[exo-4-6]{4.6}\\[2mm]
Soustraction et opposé & \hyperref[exo-4-4]{4.4}, \hyperref[exo-4-10]{4.10}\\[2mm]
Parenthèses et sommes algébriques & \hyperref[exo-4-7]{4.7}, \hyperref[exo-4-8]{4.8}, \hyperref[exo-4-9]{4.9}\\[2mm]
Variations, distances et problèmes & \hyperref[exo-4-11]{4.11}, \hyperref[exo-4-12]{4.12}, \hyperref[exo-4-13]{4.13}, \hyperref[exo-4-14]{4.14}, \hyperref[exo-4-15]{4.15}\\[2mm]
\bottomrule\end{tabularx}\par\vspace{4mm}
\otIntertitre{5. Fractions, quotients et comparaisons}\begin{tabularx}{\linewidth}{@{}Yl@{}}\toprule
Fraction et nombre quotient & \hyperref[exo-5-1]{5.1}, \hyperref[exo-5-3]{5.3}\\[2mm]
Écritures équivalentes & \hyperref[exo-5-2]{5.2}, \hyperref[exo-5-7]{5.7}, \hyperref[exo-5-9]{5.9}, \hyperref[exo-5-10]{5.10}\\[2mm]
Comparaison et droite graduée & \hyperref[exo-5-4]{5.4}, \hyperref[exo-5-5]{5.5}, \hyperref[exo-5-6]{5.6}, \hyperref[exo-5-8]{5.8}, \hyperref[exo-5-11]{5.11}, \hyperref[exo-5-12]{5.12}, \hyperref[exo-5-14]{5.14}\\[2mm]
Fractions, décimaux et pourcentages & \hyperref[exo-5-13]{5.13}, \hyperref[exo-5-15]{5.15}\\[2mm]
\bottomrule\end{tabularx}\par\vspace{4mm}
\otIntertitre{6. Additionner et soustraire des fractions}\begin{tabularx}{\linewidth}{@{}Yl@{}}\toprule
Sommes et différences de fractions & \hyperref[exo-6-1]{6.1}, \hyperref[exo-6-2]{6.2}, \hyperref[exo-6-3]{6.3}, \hyperref[exo-6-4]{6.4}\\[2mm]
Dénominateurs quelconques & \hyperref[exo-6-6]{6.6}, \hyperref[exo-6-7]{6.7}, \hyperref[exo-6-8]{6.8}, \hyperref[exo-6-9]{6.9}, \hyperref[exo-6-10]{6.10}\\[2mm]
Fraction d’une quantité & \hyperref[exo-6-5]{6.5}, \hyperref[exo-6-11]{6.11}, \hyperref[exo-6-12]{6.12}\\[2mm]
Reste et reconstitution d’un tout & \hyperref[exo-6-13]{6.13}, \hyperref[exo-6-14]{6.14}, \hyperref[exo-6-15]{6.15}\\[2mm]
\bottomrule\end{tabularx}\par\vspace{4mm}
\clearpage
\otPageTitre{Repères et exercices}{Chapitres 7 à 9}\fontsize{10}{13}\selectfont
\otIntertitre{7. Carrés et cubes}\begin{tabularx}{\linewidth}{@{}Yl@{}}\toprule
Notation du carré et du cube & \hyperref[exo-7-1]{7.1}, \hyperref[exo-7-2]{7.2}, \hyperref[exo-7-5]{7.5}\\[2mm]
Carrés de 0 à 12 et cube de 10 & \hyperref[exo-7-3]{7.3}, \hyperref[exo-7-4]{7.4}\\[2mm]
Priorités et substitution & \hyperref[exo-7-6]{7.6}, \hyperref[exo-7-7]{7.7}, \hyperref[exo-7-8]{7.8}, \hyperref[exo-7-9]{7.9}\\[2mm]
Aires, volumes et contre-exemples & \hyperref[exo-7-10]{7.10}, \hyperref[exo-7-11]{7.11}, \hyperref[exo-7-12]{7.12}, \hyperref[exo-7-13]{7.13}, \hyperref[exo-7-14]{7.14}, \hyperref[exo-7-15]{7.15}\\[2mm]
\bottomrule\end{tabularx}\par\vspace{4mm}
\otIntertitre{8. Lettres, formules et motifs}\begin{tabularx}{\linewidth}{@{}Yl@{}}\toprule
Produire et lire une formule & \hyperref[exo-8-1]{8.1}, \hyperref[exo-8-2]{8.2}, \hyperref[exo-8-4]{8.4}, \hyperref[exo-8-7]{8.7}, \hyperref[exo-8-9]{8.9}\\[2mm]
Substituer une valeur & \hyperref[exo-8-3]{8.3}, \hyperref[exo-8-6]{8.6}\\[2mm]
Suites de motifs et tableur & \hyperref[exo-8-8]{8.8}, \hyperref[exo-8-10]{8.10}, \hyperref[exo-8-11]{8.11}, \hyperref[exo-8-14]{8.14}\\[2mm]
Tester une égalité et contrôler & \hyperref[exo-8-5]{8.5}, \hyperref[exo-8-12]{8.12}, \hyperref[exo-8-13]{8.13}, \hyperref[exo-8-15]{8.15}\\[2mm]
\bottomrule\end{tabularx}\par\vspace{4mm}
\otIntertitre{9. Développer, réduire et démontrer}\begin{tabularx}{\linewidth}{@{}Yl@{}}\toprule
Développer par distributivité simple & \hyperref[exo-9-1]{9.1}, \hyperref[exo-9-2]{9.2}, \hyperref[exo-9-6]{9.6}\\[2mm]
Factoriser et réduire & \hyperref[exo-9-3]{9.3}, \hyperref[exo-9-4]{9.4}, \hyperref[exo-9-5]{9.5}, \hyperref[exo-9-7]{9.7}, \hyperref[exo-9-8]{9.8}\\[2mm]
Produire une démonstration littérale & \hyperref[exo-9-9]{9.9}, \hyperref[exo-9-11]{9.11}, \hyperref[exo-9-12]{9.12}, \hyperref[exo-9-13]{9.13}, \hyperref[exo-9-14]{9.14}\\[2mm]
Conjecture et contre-exemple & \hyperref[exo-9-10]{9.10}, \hyperref[exo-9-15]{9.15}\\[2mm]
\bottomrule\end{tabularx}\par\vspace{4mm}
\clearpage
\otPageTitre{Repères et exercices}{Chapitres 10 à 12}\fontsize{10}{13}\selectfont
\otIntertitre{10. Inconnues et équations simples}\begin{tabularx}{\linewidth}{@{}Yl@{}}\toprule
Lettre inconnue et solution & \hyperref[exo-10-4]{10.4}, \hyperref[exo-10-5]{10.5}\\[2mm]
Équations additives & \hyperref[exo-10-1]{10.1}, \hyperref[exo-10-6]{10.6}, \hyperref[exo-10-8]{10.8}\\[2mm]
Équations multiplicatives & \hyperref[exo-10-2]{10.2}, \hyperref[exo-10-3]{10.3}, \hyperref[exo-10-7]{10.7}, \hyperref[exo-10-9]{10.9}, \hyperref[exo-10-10]{10.10}\\[2mm]
Modélisation et vérification & \hyperref[exo-10-11]{10.11}, \hyperref[exo-10-12]{10.12}, \hyperref[exo-10-13]{10.13}, \hyperref[exo-10-14]{10.14}, \hyperref[exo-10-15]{10.15}\\[2mm]
\bottomrule\end{tabularx}\par\vspace{4mm}
\otIntertitre{11. Proportionnalité et pourcentages}\begin{tabularx}{\linewidth}{@{}Yl@{}}\toprule
Coefficient et procédures de proportionnalité & \hyperref[exo-11-1]{11.1}, \hyperref[exo-11-2]{11.2}, \hyperref[exo-11-5]{11.5}, \hyperref[exo-11-6]{11.6}\\[2mm]
Proportions et pourcentages & \hyperref[exo-11-3]{11.3}, \hyperref[exo-11-4]{11.4}, \hyperref[exo-11-7]{11.7}, \hyperref[exo-11-12]{11.12}, \hyperref[exo-11-13]{11.13}\\[2mm]
Prix, recettes, échelles et vitesses & \hyperref[exo-11-8]{11.8}, \hyperref[exo-11-9]{11.9}, \hyperref[exo-11-14]{11.14}, \hyperref[exo-11-15]{11.15}\\[2mm]
Tableaux et caractérisation graphique & \hyperref[exo-11-5]{11.5}, \hyperref[exo-11-10]{11.10}, \hyperref[exo-11-11]{11.11}\\[2mm]
\bottomrule\end{tabularx}\par\vspace{4mm}
\otIntertitre{12. Repérage et grandeurs en fonction d’une autre}\begin{tabularx}{\linewidth}{@{}Yl@{}}\toprule
Lire et placer des coordonnées & \hyperref[exo-12-1]{12.1}, \hyperref[exo-12-2]{12.2}, \hyperref[exo-12-3]{12.3}, \hyperref[exo-12-4]{12.4}, \hyperref[exo-12-8]{12.8}, \hyperref[exo-12-11]{12.11}\\[2mm]
Formule et tableau de valeurs & \hyperref[exo-12-6]{12.6}, \hyperref[exo-12-7]{12.7}, \hyperref[exo-12-12]{12.12}, \hyperref[exo-12-13]{12.13}, \hyperref[exo-12-15]{12.15}\\[2mm]
Courbe, nuage et dépendance & \hyperref[exo-12-5]{12.5}, \hyperref[exo-12-9]{12.9}, \hyperref[exo-12-14]{12.14}\\[2mm]
Proportionnalité graphique et limites du modèle & \hyperref[exo-12-10]{12.10}, \hyperref[exo-12-13]{12.13}, \hyperref[exo-12-15]{12.15}\\[2mm]
\bottomrule\end{tabularx}\par\vspace{4mm}
\clearpage
\otPageTitre{Repères et exercices}{Chapitres 13 à 15}\fontsize{10}{13}\selectfont
\otIntertitre{13. Angles et droites parallèles}\begin{tabularx}{\linewidth}{@{}Yl@{}}\toprule
Vocabulaire et angles usuels & \hyperref[exo-13-1]{13.1}, \hyperref[exo-13-2]{13.2}, \hyperref[exo-13-3]{13.3}, \hyperref[exo-13-4]{13.4}, \hyperref[exo-13-5]{13.5}, \hyperref[exo-13-6]{13.6}, \hyperref[exo-13-12]{13.12}\\[2mm]
Adjacents, opposés, supplémentaires, bissectrice & \hyperref[exo-13-2]{13.2}, \hyperref[exo-13-3]{13.3}, \hyperref[exo-13-4]{13.4}, \hyperref[exo-13-5]{13.5}, \hyperref[exo-13-6]{13.6}, \hyperref[exo-13-13]{13.13}\\[2mm]
Angles avec des parallèles & \hyperref[exo-13-7]{13.7}, \hyperref[exo-13-8]{13.8}, \hyperref[exo-13-11]{13.11}, \hyperref[exo-13-15]{13.15}\\[2mm]
Réciproques et raisonnement & \hyperref[exo-13-9]{13.9}, \hyperref[exo-13-10]{13.10}, \hyperref[exo-13-14]{13.14}\\[2mm]
\bottomrule\end{tabularx}\par\vspace{4mm}
\otIntertitre{14. Construire et comprendre les triangles}\begin{tabularx}{\linewidth}{@{}Yl@{}}\toprule
Angles et constructions de triangles & \hyperref[exo-14-1]{14.1}, \hyperref[exo-14-2]{14.2}, \hyperref[exo-14-3]{14.3}, \hyperref[exo-14-6]{14.6}, \hyperref[exo-14-7]{14.7}, \hyperref[exo-14-14]{14.14}\\[2mm]
Médiatrices et cercle circonscrit & \hyperref[exo-14-9]{14.9}, \hyperref[exo-14-12]{14.12}, \hyperref[exo-14-13]{14.13}\\[2mm]
Hauteurs et concours & \hyperref[exo-14-8]{14.8}, \hyperref[exo-14-12]{14.12}, \hyperref[exo-14-13]{14.13}\\[2mm]
Aires, médianes et triangles particuliers & \hyperref[exo-14-4]{14.4}, \hyperref[exo-14-5]{14.5}, \hyperref[exo-14-10]{14.10}, \hyperref[exo-14-11]{14.11}, \hyperref[exo-14-15]{14.15}\\[2mm]
\bottomrule\end{tabularx}\par\vspace{4mm}
\otIntertitre{15. Symétries et demi-tour}\begin{tabularx}{\linewidth}{@{}Yl@{}}\toprule
Symétrie axiale et construction & \hyperref[exo-15-1]{15.1}, \hyperref[exo-15-2]{15.2}, \hyperref[exo-15-6]{15.6}, \hyperref[exo-15-11]{15.11}\\[2mm]
Définition du demi-tour & \hyperref[exo-15-3]{15.3}, \hyperref[exo-15-4]{15.4}, \hyperref[exo-15-7]{15.7}, \hyperref[exo-15-8]{15.8}\\[2mm]
Conservations et images de figures & \hyperref[exo-15-5]{15.5}, \hyperref[exo-15-9]{15.9}, \hyperref[exo-15-10]{15.10}, \hyperref[exo-15-12]{15.12}\\[2mm]
Centre, milieu et justification & \hyperref[exo-15-13]{15.13}, \hyperref[exo-15-14]{15.14}, \hyperref[exo-15-15]{15.15}\\[2mm]
\bottomrule\end{tabularx}\par\vspace{4mm}
\clearpage
\otPageTitre{Repères et exercices}{Chapitres 16 à 18}\fontsize{10}{13}\selectfont
\otIntertitre{16. Parallélogrammes et aires}\begin{tabularx}{\linewidth}{@{}Yl@{}}\toprule
Définition et constructions & \hyperref[exo-16-1]{16.1}, \hyperref[exo-16-2]{16.2}, \hyperref[exo-16-6]{16.6}, \hyperref[exo-16-7]{16.7}\\[2mm]
Propriétés caractéristiques & \hyperref[exo-16-8]{16.8}, \hyperref[exo-16-9]{16.9}, \hyperref[exo-16-14]{16.14}\\[2mm]
Rectangle, losange et carré & \hyperref[exo-16-3]{16.3}, \hyperref[exo-16-5]{16.5}, \hyperref[exo-16-8]{16.8}, \hyperref[exo-16-9]{16.9}, \hyperref[exo-16-10]{16.10}\\[2mm]
Aires composées et conversions & \hyperref[exo-16-4]{16.4}, \hyperref[exo-16-11]{16.11}, \hyperref[exo-16-12]{16.12}, \hyperref[exo-16-13]{16.13}, \hyperref[exo-16-15]{16.15}\\[2mm]
\bottomrule\end{tabularx}\par\vspace{4mm}
\otIntertitre{17. Solides, disques et volumes}\begin{tabularx}{\linewidth}{@{}Yl@{}}\toprule
Perspectives, vues et dénombrement & \hyperref[exo-17-1]{17.1}, \hyperref[exo-17-11]{17.11}, \hyperref[exo-17-14]{17.14}\\[2mm]
Patrons de pavé, prisme et cylindre & \hyperref[exo-17-9]{17.9}, \hyperref[exo-17-10]{17.10}, \hyperref[exo-17-15]{17.15}\\[2mm]
Volumes et conversions de capacité & \hyperref[exo-17-2]{17.2}, \hyperref[exo-17-3]{17.3}, \hyperref[exo-17-5]{17.5}, \hyperref[exo-17-6]{17.6}, \hyperref[exo-17-12]{17.12}\\[2mm]
Aire du disque et volume du cylindre & \hyperref[exo-17-4]{17.4}, \hyperref[exo-17-7]{17.7}, \hyperref[exo-17-8]{17.8}, \hyperref[exo-17-13]{17.13}\\[2mm]
\bottomrule\end{tabularx}\par\vspace{4mm}
\otIntertitre{18. Organiser et interpréter des données}\begin{tabularx}{\linewidth}{@{}Yl@{}}\toprule
Collecte, effectifs et fréquences & \hyperref[exo-18-1]{18.1}, \hyperref[exo-18-2]{18.2}, \hyperref[exo-18-6]{18.6}, \hyperref[exo-18-12]{18.12}, \hyperref[exo-18-15]{18.15}\\[2mm]
Tableaux, diagrammes et graphiques & \hyperref[exo-18-4]{18.4}, \hyperref[exo-18-5]{18.5}, \hyperref[exo-18-7]{18.7}, \hyperref[exo-18-8]{18.8}, \hyperref[exo-18-15]{18.15}\\[2mm]
Moyenne simple et interprétation & \hyperref[exo-18-3]{18.3}, \hyperref[exo-18-9]{18.9}, \hyperref[exo-18-10]{18.10}, \hyperref[exo-18-11]{18.11}, \hyperref[exo-18-13]{18.13}\\[2mm]
Tableur et lecture critique & \hyperref[exo-18-10]{18.10}, \hyperref[exo-18-12]{18.12}, \hyperref[exo-18-13]{18.13}, \hyperref[exo-18-14]{18.14}, \hyperref[exo-18-15]{18.15}\\[2mm]
\bottomrule\end{tabularx}\par\vspace{4mm}
\clearpage
\otPageTitre{Repères et exercices}{Chapitres 19 à 20}\fontsize{10}{13}\selectfont
\otIntertitre{19. Hasard et probabilités}\begin{tabularx}{\linewidth}{@{}Yl@{}}\toprule
Expérience, issue et évènement & \hyperref[exo-19-1]{19.1}, \hyperref[exo-19-6]{19.6}, \hyperref[exo-19-9]{19.9}\\[2mm]
Échelle de probabilité et écritures & \hyperref[exo-19-2]{19.2}, \hyperref[exo-19-3]{19.3}, \hyperref[exo-19-4]{19.4}, \hyperref[exo-19-5]{19.5}\\[2mm]
Calcul en équiprobabilité & \hyperref[exo-19-6]{19.6}, \hyperref[exo-19-7]{19.7}, \hyperref[exo-19-8]{19.8}, \hyperref[exo-19-9]{19.9}, \hyperref[exo-19-12]{19.12}, \hyperref[exo-19-13]{19.13}\\[2mm]
Répétitions, effectifs et fréquences & \hyperref[exo-19-10]{19.10}, \hyperref[exo-19-11]{19.11}, \hyperref[exo-19-14]{19.14}, \hyperref[exo-19-15]{19.15}\\[2mm]
\bottomrule\end{tabularx}\par\vspace{4mm}
\otIntertitre{20. Penser et programmer par blocs}\begin{tabularx}{\linewidth}{@{}Yl@{}}\toprule
Instructions, entrées et sorties & \hyperref[exo-20-1]{20.1}, \hyperref[exo-20-5]{20.5}, \hyperref[exo-20-10]{20.10}\\[2mm]
Expressions informatiques et prévision & \hyperref[exo-20-2]{20.2}, \hyperref[exo-20-3]{20.3}, \hyperref[exo-20-4]{20.4}, \hyperref[exo-20-14]{20.14}, \hyperref[exo-20-15]{20.15}\\[2mm]
Analyse et modification de paramètres & \hyperref[exo-20-6]{20.6}, \hyperref[exo-20-10]{20.10}, \hyperref[exo-20-13]{20.13}\\[2mm]
Boucle inconditionnelle simple & \hyperref[exo-20-7]{20.7}, \hyperref[exo-20-8]{20.8}, \hyperref[exo-20-9]{20.9}, \hyperref[exo-20-11]{20.11}, \hyperref[exo-20-12]{20.12}, \hyperref[exo-20-13]{20.13}\\[2mm]
\bottomrule\end{tabularx}\par\vspace{4mm}
\clearpage
